\section{Conclusion}

We treated public-procurement question answering as typed semantic parsing followed by
deterministic execution rather than as direct factual generation.  A provenance-preserving graph
supplies the finite query universe, a closed algebra expresses set, bridge, aggregate,
comparison, and ranking computations, and the deployment runtime adds structured understanding,
conditional planning, grounding, evidence checks, bounded replanning, and controlled
non-answering.  Typed-program adaptation raises strict PACS accuracy from \PacsBaseA\% to
\PacsOursA\% and holds \PacsOursB\% under independent surface naturalisation; on a separate
system snapshot the full local runtime reaches \RuntimeOurs\% against \RuntimeBase\% untuned.

The audits matter as much as the gains.  Type validity is not semantic correctness, runtime
verification is not an oracle, exhaustive execution is exhaustive only inside a frozen snapshot,
and unresolved aliases and currencies quietly change what a correct program denotes.  The next
defensible step is not a larger headline experiment but a versioned rebuild: enforce strict
answer types and currency semantics, publish portable manifests, and measure routing,
verification, and repair under controlled ablations.  That is what would turn an executable
prototype into a reproducible account of when typed, evidence-linked question answering succeeds
and when it should decline to answer.
